\gdef\mpxshipout{\shipout\hbox\bgroup
  \setbox0=\hbox\bgroup}
\gdef\stopmpxshipout{\egroup  \dimen0=\ht0 \advance\dimen0\dp0
  \dimen1=\ht0 \dimen2=\dp0
  \setbox0=\hbox\bgroup
    \box0
    \ifnum\dimen0>0 \vrule width1sp height\dimen1 depth\dimen2 
    \else \vrule width1sp height1sp depth0sp\relax
    \fi\egroup
  \ht0=0pt \dp0=0pt \box0 \egroup}
\mpxshipout% line 31 ./p4-problem.mp
cilindro $D=0.1$%
\stopmpxshipout
\mpxshipout% line 44 ./p4-problem.mp
entrada%
\stopmpxshipout
\mpxshipout% line 45 ./p4-problem.mp
perfil parab\'olico%
\stopmpxshipout
\mpxshipout% line 46 ./p4-problem.mp
$u_{\max}=0.3$, $\bar u=0.2$%
\stopmpxshipout
\mpxshipout% line 64 ./p4-problem.mp
no deslizamiento%
\stopmpxshipout
\mpxshipout% line 66 ./p4-problem.mp
no deslizamiento%
\stopmpxshipout
\mpxshipout% line 68 ./p4-problem.mp
salida: $\nu\,\partial_n\mathbf u-p\mathbf n=0$%
\stopmpxshipout
\mpxshipout% line 76 ./p4-problem.mp
longitud $2.2$%
\stopmpxshipout
\mpxshipout% line 82 ./p4-problem.mp
altura $0.41$%
\stopmpxshipout
\mpxshipout% line 84 ./p4-problem.mp
DFG 2D-1: $\nu=0.001$, $Re=20$%
\stopmpxshipout
\mpxshipout% line 86 ./p4-problem.mp
dominio y cilindro; esquema geom\'etrico%
\stopmpxshipout
\end{document}
